% elementos.tex
%
% Copyright (C) 2022--2026 José A. Navarro Ramón <janr.devel@gmail.com>
% Licencia del código GPLv2
% Licencia Creative Commons Recognition Non-Commercial Share-alike.
% (CC-BY-NC-SA)

\chapter{Elementos en una superficie esférica}

En este texto veremos como representar distintos tipos de elementos sobre una esfera.

Para conseguirlo, se utiliza Lua\LaTeX , apoyado por el lenguaje \textsf{Lua}.
El propio formato del código fuente de este documento, puede servir para adaptarlo a
cualquier otra situación que necesite el usuario.

Utilizamos tres elementos importantes, que son los datos que se interpretan como
instrucciones para realizar una esfera y sus elementos concretos, el código fuente que
ejecuta esas instrucciones, y por último el programa \LaTeX\ que da la orden inicial para
que se obtenga el dibujo.   Se recomienda estudiar el código fuente

\begin{enumerate}
\item DATOS:

  Los datos contienen la información que define una esfera con sus elementos,
  como pueden ser meridianos, paralelos, segmentos de paralelo, arcos máximos, puntos,
  planos tangente y vectores en estos planos.
  
  Estos datos se encuentran como tablas \textsf{Lua}  en ficheros que, por comodidad se
  encuentran en un determinado directorio. En estos ejemplos es \texttt{'datoslua'}, pero
  podría ser cualquier otro que quiera el usuario.

  Es importante entender que se dibujarán todos los elementos que tengan imagen gráfica,
  y que se encuentren en el fichero de datos.
  
\item FUNCIONES:

  Las funciones Lua se encuentran en el fichero
  \texttt{lua/funciones\_esfera.lua}.
  
  La función principal es \texttt{lua/dibuja\_tikzesfera}, que según los datos
  que se le proporcionen, utilizará las funciones auxiliares que se necesiten.

\item Lua\LaTeX :

  Este será el lenguaje que organice el texto y las imágenes.
  Por ser un lenguaje general de confección de textos, se pueden utilizar los paquetes
  que se necesite, dentro del inmenso catálogo de ellos que hay a su disposición.

  El fichero principal de este documento es \texttt{esferas.tex}, que se encarga de
  organizar la estructura del documento.
  
  Para dibujar esferas con sus elementos, se utiliza el paquete \textsf{TikZ}, que se
  carga en el fichero de carga \texttt{esferas.pkg.sty} junto con otros que le pueden
  servir de complemento.

  Además, en el fichero de configuración y definiciones, \texttt{esferas.defs.sty},
  contiene, entre otros, el \LaTeX\ que se encargan de encontrar las funciones
  \textsf{Lua}, el comando \texttt{\textbackslash cargaDatos} que utiliza el fichero de
  datos que se le pasa como argumento, y \texttt{\textbackslash dibujaTikzEsfera}, que
  llama a la función \textsf{Lua} principal.

  Estos tres ficheros se encuentran en el directorio raíz.
\end{enumerate}

\section{Esferas}
Aquí nos prepararemos para representar una esfera desprovista de elementos.

\subsection{Esferas planas}
La \emph{esfera plana} se representa mediante un círculo.
El campo \texttt{'esfera'} es el que se encarga de su representación.
Este consta de los siguientes valores:

\vspace{1ex}
CAMPO \texttt{'esfera'}:
\begin{itemize}
  \setlist[itemize]{noitemsep}\vspace*{-1ex}
\item \texttt{'radio'}: Radio de la circunferencia en \unit{cm} (número).
\item \texttt{'draw'}: Indica el color que se utilizará para el borde del círculo
  (cadena de texto).
\item \texttt{'lw'}: Grosor del borde de la esfera en dimensiones \LaTeX\ '\unit{pt}'
  (cadena de texto). Ejemplo: \texttt{"0.8pt"}.
\item \texttt{'fill'}: Color de relleno del círculo (cadena de texto).
  \item \texttt{'opac'}: Opacidad del relleno (número de 0 a 1).
\end{itemize}

Lo encontramos en \texttt{'esf\_bas\_01.lua'} y lo cargamos:
Lo cargamos en Lua\LaTeX\ así:

\vspace{1ex}
\begin{tabular}{|@{\hspace{-1em}}c|}\hline
  \vphantom{\rule{0.4pt}{3ex}}
  \begin{lstlisting}[basicstyle=\ttfamily\footnotesize, language=TeX]
    \cargaDatos{datoslua/esf_bas_01.lua}
  \end{lstlisting}
  \\[1ex]
  \hline
\end{tabular}

El campo o subtabla \texttt{'esfera'} en el fichero de datos pone:

\vspace{1ex}
\begin{tabular}{|@{\hspace{-1em}}c|}\hline
  \vphantom{\rule{0.4pt}{3ex}}
  \begin{lstlisting}[basicstyle=\ttfamily\footnotesize, language=TeX]
   esfera = {
      radio=2.0, draw="black!30", lw= "0.4pt", fill="black!10", opac=1.0,
   },
  \end{lstlisting}
  \\[1ex]
  \hline
\end{tabular}

\vspace{1ex}
El código para  generar la imagen de la esfera podría ser el siguiente, aunque se
pueden utilizar otras variantes de Lua\LaTeX :

\vspace{1ex}
\begin{tabular}{|@{\hspace{-1em}}c|}\hline
  \vphantom{\rule{0.4pt}{3ex}}
  \begin{lstlisting}[basicstyle=\ttfamily\footnotesize, language=TeX]
    \cargaDatos{datoslua/esf_bas_01.lua}
    \begin{figure}[ht]
      \begin{center}
        \begin{tikzpicture}[scale=1.0]
          \dibujaTikzEsfera
        \end{tikzpicture}
      \end{center}
      \caption{Esfera básica plana, con contorno y relleno.} 
    \end{figure}
  \end{lstlisting}
  \\[1ex]
  \hline
\end{tabular}

El resultado se aprecia en la figura \ref{fig:esf-bas_plana1}

\begin{figure}[ht]
  \centering
  \cargaDatos{datoslua/esf_bas_01.lua}
  \begin{tikzpicture}[scale=1.0,]
    \dibujaTikzEsfera
  \end{tikzpicture}
  \caption{Esfera básica plana, con contorno y relleno.}
  \label{fig:esf-bas_plana1}
\end{figure}

 
Si desactivamos \texttt{fill = "--"}, conseguiremos una circunferencia sin relleno,
como se comprueba en el fichero \texttt{esf\_bas\_02.lua}:

\vspace{1ex}
\begin{tabular}{|@{\hspace{-1em}}c|}\hline
  \vphantom{\rule{0.4pt}{3ex}}
  \begin{lstlisting}[basicstyle=\ttfamily\footnotesize, language=TeX]
   esfera = {radio=2.0, draw="black!40", lw="0.8pt", fill="--" , opac=1.0,},    
  \end{lstlisting}
  \\[1ex]
  \hline
\end{tabular}

Pero, si desactivamos \texttt{draw = "--"}, como en \texttt{esf\_bas\_03.lua}, tendríamos
una esfera sin contorno, con solo relleno:

\vspace{1ex}
\begin{tabular}{|@{\hspace{-1em}}c|}\hline
  \vphantom{\rule{0.4pt}{3ex}}
  \begin{lstlisting}[basicstyle=\ttfamily\footnotesize, language=TeX]
   esfera = {radio=2.0, draw="--", lw="0.8pt", fill="black!10" , opac=1.0,},    
  \end{lstlisting}
  \\[1ex]
  \hline
\end{tabular}

\vspace{1ex}
Ambas esferas se pueden ver reunidas en la figura \ref{fig:esf-bas_plana23}:
 
\begin{figure}[ht]
  \centering
  \begin{minipage}{0.45\textwidth}
    \cargaDatos{datoslua/esf_bas_02.lua}
    \begin{tikzpicture}[scale=1.0,]
      \dibujaTikzEsfera
    \end{tikzpicture}
  \end{minipage}
  \hfill
  \begin{minipage}{0.45\textwidth}
    \cargaDatos{datoslua/esf_bas_03.lua}
    \centering
    \begin{tikzpicture}[scale=1.0,]
      \dibujaTikzEsfera
    \end{tikzpicture}
  \end{minipage}
  \caption{A la izquierda, esfera sin relleno. A la derecha sin contorno.}
  \label{fig:esf-bas_plana23}
\end{figure}


\section{Esferas con sombreado}
Para dotar a la esfera de perspectiva, se puede añadir un cierto efecto de sombreado a
través del campo \texttt{'smbresfera'}, que complementa al campo \texttt{'esfera'}, que
define el valor \texttt{'radio'} necesario para el dibujo de la esfera.
Este campo tiene los siguientes valores:

\vspace{1ex}
CAMPO \texttt{'smbresfera'}:
\begin{itemize}
  \setlist[itemize]{noitemsep}\vspace*{-1ex}
\item \texttt{'ballcolor'}: Indica el color del reflejo en la esfera.
  (cadena de texto).
  \item \texttt{'opac'}: Opacidad del relleno (número de 0 a 1).
\end{itemize}

En el fichero de datos \texttt{esf\_bas\_02.lua}, los campos \texttt{'esfera'} y \texttt{'smbresfera'} tienen:

\vspace{1ex}
\begin{tabular}{|@{\hspace{-1em}}c|}\hline
  \vphantom{\rule{0.4pt}{3ex}}
  \begin{lstlisting}[basicstyle=\ttfamily\footnotesize, language=TeX]
   esfera = {radio=2.0, draw="white", lw="0pt", fill="white", opac=0,},
   smbresfera = {ballcolor="green!50", opac=0.35},
  \end{lstlisting}
  \\[1ex]
  \hline
\end{tabular}

\vspace{1ex}
El código para  generar la imagen de la esfera podría ser el siguiente, aunque se
pueden utilizar otras variantes de Lua\LaTeX :

\vspace{1ex}
\begin{tabular}{|@{\hspace{-1em}}c|}\hline
  \vphantom{\rule{0.4pt}{3ex}}
  \begin{lstlisting}[basicstyle=\ttfamily\footnotesize, language=TeX]
    \cargaDatos{datoslua/esf_bas_02.lua}
    \begin{figure}[ht]
      \begin{center}
        \begin{tikzpicture}[scale=1.0]
          \dibujaTikzEsfera
        \end{tikzpicture}
      \end{center}
      \caption{Esfera básica plana.} 
    \end{figure}
  \end{lstlisting}
  \\[1ex]
  \hline
\end{tabular}

\vspace{1ex}
Como resultado,
%figura \ref{fig:esf-sombr_verde},
se genera una superficie esférica
limpia, pues el campo \texttt{'esfera'} es el único en el fichero de datos.

\vspace{1ex}
Si cambiamos los valores de los campos, como en \texttt{esf\_bas\_03.lua}:

\vspace{1ex}
\begin{tabular}{|@{\hspace{-1em}}c|}\hline
  \vphantom{\rule{0.4pt}{3ex}}
  \begin{lstlisting}[basicstyle=\ttfamily\footnotesize, language=TeX]
   esfera = {radio=2.0, draw="white", lw= "0pt", fill="white", opac=0,},
   smbresfera = {ballcolor="white", opac=0.35},
  \end{lstlisting}
  \\[1ex]
  \hline
\end{tabular}


\begin{figure}[ht]
  \centering
  \begin{minipage}{0.45\textwidth}
    \cargaDatos{datoslua/esf_smbr_01.lua}
    \begin{tikzpicture}[scale=1.0,]
      \dibujaTikzEsfera
    \end{tikzpicture}
  \end{minipage}
  \hfill
  \begin{minipage}{0.45\textwidth}
    \cargaDatos{datoslua/esf_smbr_02.lua}
    \centering
    \begin{tikzpicture}[scale=1.0,]
      \dibujaTikzEsfera
    \end{tikzpicture}
  \end{minipage}
  \caption{Esferas sombreadas que muestran una cierta perspectiva, en verde y en gris.}
  \label{fig:esf-sombreadas}
\end{figure}



\section{Meridianos}
Para representar uno o más meridianos en la esfera, tenemos que utilizar el comando:
\begin{center}
  {\small\texttt{\textbackslash dibujaTikZEsfera[true o false]}}
\end{center}
Para representar los meridianos, el comando Lua\LaTeX\ anterior necesita saber si la
esfera es transparente o no. Esto es necesario, pues en caso de que lo sea, se pueden
ver los meridianos que serían tapados por la esfera.

Por tanto, hay que indicar, entre corchetes, \texttt{true} si es transparente o
\texttt{false}, si es opaca.
Si no se detallara ningún argumento, el valor por defecto sería \texttt{false} lo que
indicaría que la esfera es opaca, por lo que no dejaría ver los elementos que se
encuentran tapados.

Con el comando anterior se leen, además de \texttt{esfera}, oros campos o subtablas,
como \texttt{'observador'}, \texttt{'meridprops'} y \texttt{'meridianos'} para representar
meridianos, pues su representación dependería del ángulo de observación de la esfera:

\begin{enumerate}
\item El campo \texttt{'observador'} define la vista de un observador situado muy lejos
  de la esfera, definida por las coordenadas angulares esféricas polar $\theta$ y
  azimutal $\phi$ en grados:
  \begin{itemize}
    \setlist[itemize]{noitemsep}\vspace*{-1ex}
  \item \texttt{thetaD}: Ángulo polar de observación de la esfera, en grados.
  \item \texttt{phiD}: Ángulo azimutal de observación, en grados.
  \end{itemize}
\item El campo \texttt{'meridprops'} define propiedades visuales por defecto del dibujo
  de los meridianos, y otras de la parte \emph{visible} y la \emph{invisible} de los
  meridianos (si la esfera fuera opaca):
  \begin{itemize}
    \setlist[itemize]{noitemsep}\vspace*{-1ex}
  \item \texttt{loops}: Cuántas veces se repite el bucle de trazado de cada meridano
    (de cuantos puntos consta cada meridiano).
  \item \texttt{color\_vis}: Color reservado para el dibujo de la parte \emph{visible}.
  \item \texttt{lw\_vis}: Grosor de línea de la parte \emph{visible}.
  \item \texttt{color\_novis}: Color reservado para el dibujo de la parte
    \emph{invisible}.
  \item \texttt{lw\_novis}: Grosor de línea de la parte \emph{invisible} (dimensión
    \LaTeX , preferiblemente en puntos (\unit{pt}))
  \end{itemize}  
\item El campo \texttt{'meridianos'}, define el ángulo azimutal $\phi$ que define la
  longitud en grados del meridiano:
  \begin{itemize}
        \setlist[itemize]{noitemsep}\vspace*{-1ex}
        \item \texttt{phiD}: Ángulo azimutal que define el meridiano, en grados (número)
        \end{itemize}
        Además del valor anterior, se puede destacar más o menos uno o más meridianos
        con respecto a los demás, añadiendo todos o alguno de los valores del
        campo \texttt{'meridprops'}. Ya veremos esto en algún ejemplo.
\end{enumerate}

Veamos uno por uno los campos en los ficheros que mostraremos en los siguientes
ejemplos, \texttt{esf\_merid\_01.lua} y \texttt{esf\_merid\_02.lua}:

\vspace{1ex}
CAMPO \texttt{'observador'}:

\begin{tabular}{|@{\hspace{-0.8em}}c|}\hline
  \vphantom{\rule{0.4pt}{3ex}}
  \begin{lstlisting}[style=myLuastyle]
    observador = {thetaD = 65, phiD = 15},
  \end{lstlisting}
  \\[1ex]
  \hline
\end{tabular}

\vspace{1ex}
CAMPO \texttt{'meridprops'}:

\begin{tabular}{|@{\hspace{-0.8em}}c|}\hline
  \vphantom{\rule{0.4pt}{3ex}}
  \begin{lstlisting}[style=myLuastyle]
   meridprops = {
      loops = 120, color_vis = "black!60", lw_vis="0.4pt",
      color_novis = "black!30", lw_novis="0.3pt",
   },
  \end{lstlisting}
  \\[1ex]
  \hline
\end{tabular}

\vspace{1ex}
CAMPO \texttt{'meridianos'}:

\begin{tabular}{|@{\hspace{-0.8em}}c|}\hline
  \vphantom{\rule{0.4pt}{3ex}}
  \begin{lstlisting}[style=myLuastyle]
   meridianos = {
     {phiD = 0,}, {phiD = 30}, {phiD = 60},
     {phiD = 90}, {phiD = 120}, {phiD = 150},
     {phiD = 180}, {phiD = -30}, {phiD = -60},
     {phiD = -90}, {phiD = -120},{phiD = -150}
   },
  \end{lstlisting}
  \\[1ex]
  \hline
\end{tabular}


En el fichero \texttt{esf\_merid\_01.lua} se define una esfera plana:

\vspace{1ex}
\begin{tabular}{|@{\hspace{-1em}}c|}\hline
  \vphantom{\rule{0.4pt}{3ex}}
  \begin{lstlisting}[basicstyle=\ttfamily\footnotesize, language=TeX]
   esfera = {radio=2.0, draw="black!25", lw="0.6pt", fill="black!8", opac=1.0,},
  \end{lstlisting}
  \\[1ex]
  \hline
\end{tabular}

La rutina de trazado de meridianos se encarga de dibujar primero la parte
\emph{invisible} (en caso de que la esfera sea translúcida), y solo después,
la parte \emph{visible}.

Para no repetir el mismo código, resaltaremos el fichero de datos y el comando
\LaTeX\ que utilizamos Para una esfera opaca sería:

\vspace{1ex}
\begin{tabular}{|@{\hspace{-1em}}c|}\hline
  \vphantom{\rule{0.4pt}{3ex}}
  \begin{lstlisting}[basicstyle=\ttfamily\footnotesize, language=TeX]
    \cargaDatos{datoslua/esf_merid_01.lua}
        ...   
        \dibujaTikzEsfera
        ...
  \end{lstlisting}
  \\[1ex]
  \hline
\end{tabular}

El fichero \texttt{esf\_merid\_02.lua}, define una esfera sombreada (el resto de
campos es el mismo que en el fichero \texttt{esf\_merid\_01.lua}:

\vspace{1ex}
\begin{tabular}{|@{\hspace{-1em}}c|}\hline
  \vphantom{\rule{0.4pt}{3ex}}
  \begin{lstlisting}[basicstyle=\ttfamily\footnotesize, language=TeX]
   esfera = {radio=2.0, draw="black!25", lw="0.6pt", fill="black!8", opac=1.0,},
   smbresfera = {color="gray!1", opac=0.2},
  \end{lstlisting}
  \\[1ex]
  \hline
\end{tabular}

A continuación se muestran las dos esferas:

 \begin{figure}[ht]
  \centering
  \begin{minipage}{0.45\textwidth}
    \cargaDatos{datoslua/esf_merid_01.lua}
    \begin{tikzpicture}[scale=1.0,]
      \dibujaTikzEsfera
    \end{tikzpicture}
    \caption{Esfera plana.}
  \end{minipage}
  \hfill
  \begin{minipage}{0.45\linewidth}
    \cargaDatos{datoslua/esf_merid_02.lua}
     \centering
     \begin{tikzpicture}[scale=1.0,]
       \dibujaTikzEsfera
     \end{tikzpicture}
     \caption{Esfera sombreada.}
   \end{minipage}
   \caption{Esferas opacas.}
 \end{figure}


Ahora dibujamos las mismas esferas y meridianos, pero considerando que son translúcidas
(añadiendo \texttt{true} al comando: \texttt{\textbackslash dibujaTikzEsfera[true]}

\begin{figure}[ht]
  \centering
  \begin{minipage}[b]{0.45\textwidth}
    \cargaDatos{datoslua/esf_merid_01.lua}
    \begin{tikzpicture}[scale=1.0,]
      \dibujaTikzEsfera[true]
    \end{tikzpicture}
    \caption{Esfera plana.}
  \end{minipage}
  \hfill
  \begin{minipage}[b]{0.45\linewidth}
    \cargaDatos{datoslua/esf_merid_02.lua}
    \centering
    \begin{tikzpicture}[scale=1.0,]
      \dibujaTikzEsfera[true]
    \end{tikzpicture}
    \caption{Esfera sombreada.}
  \end{minipage}
  \caption{Esferas translúcidas.}
\end{figure}

Otras figuras:

En la siguientes figuras opacas resaltamos un meridiano.
Para ello, utilizamos el fichero \texttt{esf\_meridres\_01.lua} para la esfera
gris:

\vspace{1ex}
\begin{tabular}{|@{\hspace{-1em}}c|}\hline
  \vphantom{\rule{0.4pt}{3ex}}
  \begin{lstlisting}[basicstyle=\ttfamily\footnotesize, language=TeX]
    meridprops = {
      loops = 120, color_vis = "black!60", lw_vis="0.4pt",
      color_novis = "black!30", lw_novis="0.3pt",
    },
  \end{lstlisting}
  \\[1ex]
  \hline
\end{tabular}

Se puede observar que para esta esfera se resaltan de rojo los puntos visibles y
naranja los invisibles, los meridianos \qty{0}{\degree} y \qty{180}{\degree}:

\vspace{1ex}
\begin{tabular}{|@{\hspace{-1em}}c|}\hline
  \vphantom{\rule{0.4pt}{3ex}}
  \begin{lstlisting}[basicstyle=\ttfamily\footnotesize, language=TeX]
    meridianos = {
      {phiD = 0,
        color_vis="red", lw_vis="0.8pt", color_novis="orange", lw_novis="0.8pt"},
      {phiD = 30}, {phiD = 60}, {phiD = 90}, {phiD = 120}, {phiD = 150},
      {phiD = 180,
        color_vis="red", lw_vis="0.8pt", color_novis="orange", lw_novis="0.8pt"},
      {phiD = -30}, {phiD = -60}, {phiD = -90},  {phiD = -120}, {phiD = -150}
    },
  \end{lstlisting}
  \\[1ex]
  \hline
\end{tabular}

Para la verde utilizamos \texttt{esf\_meridres\_02.lua}:

\vspace{1ex}
\begin{tabular}{|@{\hspace{-1em}}c|}\hline
  \vphantom{\rule{0.4pt}{3ex}}
  \begin{lstlisting}[basicstyle=\ttfamily\footnotesize, language=TeX]
    meridprops = {
      loops = 120, color_vis = "black!80", lw_vis="0.6pt",
      color_novis = "black!40", lw_novis="0.5pt",
    },
  \end{lstlisting}
  \\[1ex]
  \hline
\end{tabular}

Para la esfera verde se resaltan de rojo los puntos visibles y en  naranja los
invisibles, los meridianos \qty{-30}{\degree} y \qty{150}{\degree}:

\vspace{1ex}
\begin{tabular}{|@{\hspace{-1em}}c|}\hline
  \vphantom{\rule{0.4pt}{3ex}}
  \begin{lstlisting}[basicstyle=\ttfamily\footnotesize, language=TeX]
    meridianos = {
      {phiD = 0,}, {phiD = 30}, {phiD = 60}, {phiD = 90}, {phiD = 120},
      {phiD = 150,
        color_vis="red", lw_vis="1pt", color_novis="orange", lw_novis="1pt"},
      {phiD = 180},
      {phiD = -30,
        color_vis="red", lw_vis="1pt", color_novis="orange", lw_novis="1pt"},
      {phiD = -60}, {phiD = -90}, {phiD = -120}, {phiD = -150}
    },
  \end{lstlisting}
  \\[1ex]
  \hline
\end{tabular}

\begin{figure}[ht]
  \centering
  \begin{minipage}[b]{0.45\textwidth}
    \cargaDatos{datoslua/esf_meridres_01.lua}
    \begin{tikzpicture}[scale=1.0,]
      \dibujaTikzEsfera
    \end{tikzpicture}
    \caption{Esfera gris.}
  \end{minipage}
  \hfill
  \begin{minipage}[b]{0.45\linewidth}
    \cargaDatos{datoslua/esf_meridres_02.lua}
    \centering
    \begin{tikzpicture}[scale=1.0,]
      \dibujaTikzEsfera
    \end{tikzpicture}
    \caption{Esfera verde.}
  \end{minipage}
  \caption{Esferas opacas con dos meridianos resaltados.}
\end{figure}


\begin{figure}[ht]
  \centering
  \begin{minipage}[b]{0.45\textwidth}
    \cargaDatos{datoslua/esf_meridres_01.lua}
    \begin{tikzpicture}[scale=1.0,]
      \dibujaTikzEsfera[true]
    \end{tikzpicture}
    \caption{Esfera gris.}
  \end{minipage}
  \hfill
  \begin{minipage}[b]{0.45\linewidth}
    \cargaDatos{datoslua/esf_meridres_02.lua}
    \centering
    \begin{tikzpicture}[scale=1.0,]
      \dibujaTikzEsfera[true]
    \end{tikzpicture}
    \caption{Esfera verde.}
  \end{minipage}
  \caption{Esferas translúcidas con dos meridianos resaltados.}
\end{figure}


\subsection{Paralelos y polos}
En este apartado veremos dos elementos que se complementan, por lo que aparecen juntos,
aunque no es necesario. Se trata de los paralelos de una esfera y sus los polos norte
y sur (los puntos con $\theta=\qty{0}{\degree}$ y $\theta=\qty{180}{\degree}$,
respectivamente.

Como en todos los casos, para representar uno o más polos y/o paralelos en la esfera, se
utiliza el comando principal definido para Lua\LaTeX :
\begin{center}
  {\small\texttt{\textbackslash dibujaTikZEsfera[true]}}
  
  {\small\texttt{\textbackslash dibujaTikZEsfera[false]}}
\end{center}

Recordamos que, con \texttt{true} se indica que la esfera es total o parcialmente
transparente (translúcida) y con \texttt{false}, que es opaca.
De la misma manera que con los meridianos, si no se indica nada, se supone que es opaca:
\begin{center}
  {\small\texttt{\textbackslash dibujaTikZEsfera}}
\end{center}

Con el comando anterior se leen, además de \texttt{esfera}, oros campos o subtablas,
como \texttt{'observador'}, \texttt{'paralprops'} y \texttt{'paralelos} para representar
paralelos. Todo esto, excepto \texttt{'paralelos'} es igual que con los meridianos,
aunque lo repetiremos aquí:

%\begin{enumerate}
%\item El campo \texttt{'observador'} define la vista de un observador situado muy lejos
%  de la esfera, definida por las coordenadas angulares esféricas polar $\theta$ y
%  azimutal $\phi$ en grados:
%  \begin{itemize}
%    \setlist[itemize]{noitemsep}\vspace*{-1ex}
%  \item \texttt{thetaD}: Ángulo polar de observación de la esfera, en grados.
%  \item \texttt{phiD}: Ángulo azimutal de observación, en grados.
%  \end{itemize}
%\item El campo \texttt{'paralprops'} define propiedades visuales por defecto del dibujo
%  de los meridianos, y otras de la parte \emph{visible} y la \emph{invisible} de los
%  meridianos (si la esfera fuera opaca):
%  \begin{itemize}
%    \setlist[itemize]{noitemsep}\vspace*{-1ex}
%  \item \texttt{loops}: Cuántas veces se repite el bucle de trazado de cada meridano
%    (de cuantos puntos consta cada meridiano).
%  \item \texttt{color\_vis}: Color reservado para el dibujo de la parte \emph{visible}.
%  \item \texttt{lw\_vis}: Grosor de línea de la parte \emph{visible}.
%  \item \texttt{color\_novis}: Color reservado para el dibujo de la parte
%    \emph{invisible}.
%  \item \texttt{lw\_novis}: Grosor de línea de la parte \emph{invisible} (dimensión
%    \LaTeX , preferiblemente en puntos (\unit{pt}))
%  \end{itemize}  
%\item El campo \texttt{'paralelos'}, define el ángulo polar $\theta$ que define el
%  paralelo en grados:
%  \begin{itemize}
%    \setlist[itemize]{noitemsep}\vspace*{-1ex}
%  \item \texttt{thetaD}: Ángulo azimutal que define el meridiano, en grados (número)
%  \end{itemize}
%  Además del valor anterior, se puede destacar más o menos uno o más paralelos
%  con respecto a los demás, añadiendo todos o alguno de los valores del
%  campo \texttt{'paralprops'}, excepto \texttt{'loops'}.
%\item El campo \texttt{'polprops'} define las características de los puntos que
%  representan los polos norte y sur de la esfera:
%  \begin{itemize}
%  \item \texttt{'color\_vis'}: Color de los polos visibles (cadena de texto).
%  \item \texttt{'radio\_vis'}: Radio en dimensiones \LaTeX , preferiblemente \unit{pt}
%    de los polos visibles (cadena de texto).
%  \item \texttt{'color\_novis'}: Color de los polos invisibles (cadena de texto).
%  \item \texttt{'radio\_novis'}: Radio en dimensiones \LaTeX , preferiblemente \unit{pt}
%    de los polos invisibles (cadena de texto). 
%  \end{itemize}
%\end{enumerate}
%
%Veamos uno por uno los campos en los ficheros que mostraremos en los siguientes
%ejemplos, \texttt{esf\_paral\_01.lua} y \texttt{esf\_paral\_02.lua}, obviando
%los campos
%\begin{center}\texttt{'esfera'}, \texttt{smbresfera} y \texttt{'observador'}\end{center}
%
%\vspace{1ex}
%CAMPO \texttt{'polprops'}:
%
%\begin{tabular}{|@{\hspace{-0.8em}}c|}\hline
%  \vphantom{\rule{0.4pt}{3ex}}
%  \begin{lstlisting}[style=myLuastyle]
%    polprops = {
%      color_vis = "black!70", radio_vis="1pt",
%      color_novis = "black!40", radio_novis="1pt",
%    },
%  \end{lstlisting}
%  \\[1ex]
%  \hline
%\end{tabular}
%
%
%\vspace{1ex}
%CAMPO \texttt{'paralprops'}:
%
%\begin{tabular}{|@{\hspace{-0.8em}}c|}\hline
%  \vphantom{\rule{0.4pt}{3ex}}
%  \begin{lstlisting}[style=myLuastyle]
%    paralprops = {
%      loops = 240, color_vis = "black!60", lw_vis="0.4pt",
%      color_novis = "black!30", lw_novis="0.3pt",
%    },
%  \end{lstlisting}
%  \\[1ex]
%  \hline
%\end{tabular}
%
%\vspace{1ex}
%CAMPO \texttt{'paralelos'}:
%
%\begin{tabular}{|@{\hspace{-0.8em}}c|}\hline
%  \vphantom{\rule{0.4pt}{3ex}}
%  \begin{lstlisting}[style=myLuastyle]
%   meridianos = {
%     {phiD = 30}, {phiD = 60},
%     {phiD = 90}, {phiD = 120}, {phiD = 150},
%     {phiD = 180}, {phiD = -30}, {phiD = -60},
%     {phiD = -90}, {phiD = -120},{phiD = -150}
%   },
%  \end{lstlisting}
%  \\[1ex]
%  \hline
%\end{tabular}
%
%
%En el fichero \texttt{esf\_paral\_01.lua} se define una esfera plana:
%
%\vspace{1ex}
%\begin{tabular}{|@{\hspace{-1em}}c|}\hline
%  \vphantom{\rule{0.4pt}{3ex}}
%  \begin{lstlisting}[basicstyle=\ttfamily\footnotesize, language=TeX]
%    esfera = {radio=2.0, draw="black!25", lw="0.6pt", fill="black!8", opac=1.0,},
%  \end{lstlisting}
%  \\[1ex]
%  \hline
%\end{tabular}
%
%Mientras que en \texttt{esf\_paral\_02.lua} está sombreada
%
%\vspace{1ex}
%\begin{tabular}{|@{\hspace{-1em}}c|}\hline
%  \vphantom{\rule{0.4pt}{3ex}}
%  \begin{lstlisting}[basicstyle=\ttfamily\footnotesize, language=TeX]
%   esfera = {radio=2.0, draw="black!25", lw="0.6pt", fill="black!8", opac=1.0,},
%   smbresfera = {ballcolor="white", opac=0.4},
%  \end{lstlisting}
%  \\[1ex]
%  \hline
%\end{tabular}
%
%\begin{figure}[ht]
%  \centering
%  \begin{minipage}{0.45\textwidth}
%    \cargaDatos{datoslua/esf_paral_01.lua}
%    \begin{tikzpicture}[scale=1.0,]
%      \dibujaTikzEsfera
%    \end{tikzpicture}
%  \end{minipage}
%  \hfill
%  \begin{minipage}{0.45\linewidth}
%    \cargaDatos{datoslua/esf_paral_02.lua}
%    \centering
%    \begin{tikzpicture}[scale=1.0,]
%      \dibujaTikzEsfera
%    \end{tikzpicture}
%  \end{minipage}
%  \caption{Esferas opacas, plana y sombreada.}
%\end{figure}
%
%\begin{figure}[ht]
%  \centering
%  \begin{minipage}{0.45\textwidth}
%    \cargaDatos{datoslua/esf_paral_01.lua}
%    \begin{tikzpicture}[scale=1.0,]
%      \dibujaTikzEsfera[true]
%    \end{tikzpicture}
%  \end{minipage}
%  \hfill
%  \begin{minipage}{0.45\linewidth}
%    \cargaDatos{datoslua/esf_paral_02.lua}
%    \centering
%    \begin{tikzpicture}[scale=1.0,]
%      \dibujaTikzEsfera[true]
%    \end{tikzpicture}
%  \end{minipage}
%  \caption{Esferas transparentes, plana y sombreada.}
%\end{figure}


%\vspace{1ex}
%En el fichero \texttt{esf\_merpar\_01.lua} representamos meridianos y paralelos
%en la esfera:
%
%\begin{figure}[ht]
%  \centering
%  \begin{minipage}{0.45\textwidth}
%    \cargaDatos{datoslua/esf_merpar_01.lua}
%    \begin{tikzpicture}[scale=1.0,]
%      \dibujaTikzEsfera
%    \end{tikzpicture}
%  \end{minipage}
%  \hfill
%  \begin{minipage}{0.45\linewidth}
%    \cargaDatos{datoslua/esf_merpar_01.lua}
%    \centering
%    \begin{tikzpicture}[scale=1.0,]
%      \dibujaTikzEsfera[true]
%    \end{tikzpicture}
%  \end{minipage}
%  \caption{Esferas opaca y translúcida con varios meridianos y paralelos.}
%\end{figure}


%\subsection{Segmentos de paralelos}
%
%\begin{figure}[ht]
%  \centering
%  \begin{minipage}{0.45\textwidth}
%    \cargaDatos{datoslua/esf_sgmpar_01.lua}
%    \begin{tikzpicture}[scale=1.0,]
%      \dibujaTikzEsfera
%    \end{tikzpicture}
%  \end{minipage}
%  \hfill
%  \begin{minipage}{0.45\linewidth}
%    \cargaDatos{datoslua/esf_sgmpar_02.lua}
%    \centering
%    \begin{tikzpicture}[scale=1.0,]
%      \dibujaTikzEsfera[true]
%    \end{tikzpicture}
%  \end{minipage}
%    \caption{Esferas opacas con un segmento paralelo. La de la izquierda es plana.
%      La de la derecha está sombreada.}
%\end{figure}
%
%
%
%Hola.


\end{document}











%%% Local Variables:
%%% coding: utf-8
%%% mode: latex
%%% TeX-engine: luatex
%%% TeX-master: "../esferas.tex"
%%% End:


